% =============================================================
% File: chapters/ta/chapter-3/rancangan/rancangan-arsitektur-sistem.tex
%
% PERUBAHAN:
%   - "Audit Trail Service (ATS)" → "Audit Log Service (ALS)"
%   - "ATS" → "ALS"
%   - Caption figure: "Audit Trail" → "Audit Log"
%   - "rekam jejak aktivitas" → "aktivitas"
%   - "immutable" pada klaim lokal → "tamper-evident"
%   - Paragraf satu baris
% =============================================================

\subsection{Rancangan Arsitektur Solusi}

Guna mengatasi permasalahan terfragmentasinya pencatatan dan kerentanan integritas log pada DecMed, penelitian ini mengusulkan penambahan satu komponen utama pada arsitektur sistem, yaitu \textbf{Audit Log Service (ALS)}. ALS dirancang sebagai layanan \textit{backend} independen yang menjembatani pengumpulan bukti digital dari seluruh komponen DecMed menuju media penyimpanan terdesentralisasi.

Penambahan ALS tidak mengubah alur bisnis maupun fungsionalitas utama DecMed yang telah berjalan pada komponen eksisting, seperti \textit{Patient Client}, \textit{Hospital Client}, PRE Server, maupun Gas Station Server. ALS bertindak sebagai pengumpul log terpusat yang memanfaatkan kembali (\textit{reuse}) infrastruktur \textit{off-chain} (IPFS Cluster) dan \textit{on-chain} (jaringan IOTA) yang sudah tersedia pada DecMed.

Gambaran arsitektur solusi sistem DecMed beserta integrasi ALS ditunjukkan pada Gambar~\ref{fig:arsitektur-solusi}. Warna merah merepresentasikan komponen baru dan penyesuaian alur interaksi yang ditambahkan pada arsitektur DecMed.

\begin{figure}[htbp]
\centering
\includegraphics[width=1\textwidth]{chapter-3/decmed_ats_architecture.png}
\caption{Rancangan Arsitektur Solusi Audit Log DecMed}
\label{fig:arsitektur-solusi}
\end{figure}

Berdasarkan Gambar~\ref{fig:arsitektur-solusi}, alur interaksi dan integrasi ALS dalam arsitektur DecMed dirancang sebagai berikut.

\begin{enumerate}
    \item \textbf{Pengiriman Event dari Komponen DecMed:} Setiap kali terjadi aktivitas bernilai audit pada aplikasi klien (\textit{Patient Client} dan \textit{Hospital Client}) maupun komponen perantara (PRE Server dan Gas Station Server), komponen tersebut mengirimkan rekaman \textit{audit event} ke ALS melalui antarmuka komunikasi yang aman.
    \item \textbf{Pemrosesan dan Pengelompokan Log:} ALS menerima \textit{event} tersebut, menyelaraskan formatnya sesuai skema terstruktur, membentuk rantai integritas (\textit{hash chain}), dan mengelompokkan rekaman log ke dalam satu berkas arsip (\textit{batch}) secara periodik.
    \item \textbf{Penyimpanan Berkas \textit{Off-Chain} di IPFS:} ALS mengunggah berkas arsip log ke IPFS Cluster eksisting DecMed dan menerapkan perintah \textit{pinning} agar berkas log tidak terhapus oleh proses \textit{Garbage Collection}.
    \item \textbf{Pencatatan Metadata \textit{On-Chain} di IOTA:} ALS menerbitkan transaksi ke jaringan IOTA melalui Gas Station Server untuk mencatatkan metadata berkas ke dalam \textit{smart contract} Move sebagai \textit{immutable trust anchor}.
\end{enumerate}

Dengan rancangan arsitektur ini, ALS berhasil memisahkan beban pemrosesan dan penyimpanan log dari transaksi data klinis utama, sekaligus menjamin bahwa seluruh aktivitas pada DecMed terekam secara utuh, bersifat \textit{tamper-evident} pada level lokal, dan \textit{immutable} setelah tersimpan pada IOTA dan IPFS.